\section{At the foot of the altar (nn.~1013--1023)}

The 1962 \work{Ordo Missae} opens with the priest already vested and already
at the altar step, having made the reverence due to the altar. He signs
himself and says aloud, unless a particular rubric provides otherwise, the
trinitarian formula; then, hands joined before his breast, he begins an
antiphon.

\begin{appointed}{1013--1014 $\cdot$ \latincue{In nomine Patris} and the antiphon}
\textit{Focused source locator.} The appointed unit is identified by this incipit and marginal number; only phrases required by the analysis below are quoted.
\end{appointed}

\subsection*{The psalm and its condition}

What follows is Psalm 42, verses 1 to 5, said alternately with the ministers,
with \latincue{Gloria Patri} and the antiphon repeated. The missal prints the
alternation with the letters \latincue{S.} and \latincue{M.} --- \latincue{sacerdos}
and \latincue{minister} --- and it prints the psalm reference in the margin,
\latincue{Ps.~42, 1-5}, in the Vulgate numbering that the Douay-Rheims Bible
also uses.

\begin{appointed}{1015 $\cdot$ \latincue{Iudica me, Deus} (Ps.~42, 1--5)}
\textit{Focused source locator.} The appointed unit is identified by this incipit and marginal number; only phrases required by the analysis below are quoted.
\end{appointed}

\begin{rubricnote}{Rubric printed at n.~1016, and \work{Rubricae generales Missalis romani} nn.~424--425}
In Requiem Masses, and in Masses of the season from Passion Sunday to Maundy
Thursday inclusive, the psalm \latincue{Iudica me, Deus} is omitted together
with \latincue{Gloria Patri} and the repetition of the antiphon. The priest
says \latincue{In nomine Patris}, \latincue{Introibo} and
\latincue{Adiutorium}, and then makes the confession. The omission is
therefore not a general Lenten rule: it begins at Passion Sunday, and n.~425
restricts it to Masses \latincue{de Tempore}, not to a saint's Mass falling in
that fortnight.\par
A second and wider omission is governed by n.~424, which prescribes the psalm
with its antiphon and the confession with the absolution before the altar
steps in every Mass, sung or read, but then removes the whole block ---
psalm, antiphon, confession, absolution, the following versicles, and both
\latincue{Aufer a nobis} and \latincue{Oramus te} --- from the Masses that
follow the blessing of candles at the Purification, the blessing and imposition
of ashes, the blessing and procession of palms, the Easter Vigil, the Rogation
procession, and certain Masses following consecrations in the \work{Pontificale
Romanum}. The principle is legible: where the people have already been
gathered, blessed and led in procession, the priest's private approach to the
altar has already happened in public, and the rite does not repeat it.
\end{rubricnote}

The psalm is the reason the rite has its shape. Its speaker is a man barred
from the sanctuary and asking to be let in: \latincue{discerne causam meam de
gente non sancta}, distinguish my cause from a nation that is not holy. Its
fourth verse supplies the antiphon, so that the antiphon is a quotation from
the psalm's own centre and the whole psalm is heard as an expansion of the one
line the priest has already said twice. The Douay-Rheims renders that verse
``And I will go in to the altar of God: to God who giveth joy to my youth''
(Ps.~42, 4). The Latin \latincue{iuventutem meam} carries no suggestion that
the speaker is young; it is the joy that renews, and the older English witness
prints both ``who giveth joy to my youth'' at the antiphon's first appearance
and ``who rejoiceth my youth'' at its repetition, which shows how little the
translator thought turned on the phrase.

The omission of the psalm in Passiontide and at Requiems is the rite's own
commentary on it. The psalm asks for light and truth to lead the speaker to the
holy mountain; from Passion Sunday the Church has already begun to keep silence
about that leading, and at a Requiem the light asked for is asked in the
introit instead. The omission is not a shortening for convenience --- the
\latincue{Gloria Patri} disappears with the psalm at exactly the same seasons
at which it disappears from the \latincue{Lavabo} later in the same Mass.

\subsection*{The confession}

\begin{appointed}{1016--1017 $\cdot$ \latincue{Adiutorium} and the \latincue{Confiteor}}
\textit{Focused source locator.} The appointed unit is identified by this incipit and marginal number; only phrases required by the analysis below are quoted.
\end{appointed}

\begin{englishwitness}{xvi--xvii}
``I confess to Almighty God, to blessed Mary, ever a Virgin, to blessed Michael
the Archangel, to blessed John Baptist, to the holy Apostles Peter and Paul, to
all the saints, and to you, Father, that I have sinned exceedingly in thought,
word and deed, \emph{through my fault, through my fault, through my most
grievous fault}. Therefore I beseech the Blessed Mary, ever a virgin, Blessed
Michael the Archangel, Blessed John Baptist, the holy Apostles Peter and Paul,
and all the saints, and you, O Father, to pray to the Lord our God for me.''
The 1861 book prints the servers' form in full and abbreviates the priest's to
\latincue{Confiteor Deo omnipotenti, \&c.}, because in a lay hand missal the
form the reader hears twice is the one worth printing whole.
\end{englishwitness}

The \work{Confiteor} as the 1962 book prints it is a fixed list: God, then the
Virgin, then Michael, John the Baptist, Peter and Paul, then all the saints,
and then the person present. The missal prints no variant and no invitation to
add a name, and the general rubrics that govern this part of the Mass are
concerned only with when the confession is said at all, not with what it may
contain. That silence is itself the point. Preparation of this kind had been
the freest part of the medieval Mass, and by 1962 it had become the most
tightly fixed. Where
the Sarum use had the celebrant say \latincue{Veni Creator} while vesting,
Psalm 42 on the way to the altar, and a short \work{Confiteor} at the step,
and where the Dominicans, Carthusians and Carmelites said the antiphon without
the psalm at all, the Roman book of 1570 fixed one form and the code of 1960
sealed it against local additions.

The threefold \latincue{mea culpa} with the striking of the breast is printed
as a rubric inside the prayer, not before it: the missal will not let the
gesture float free of the words it accompanies. The same technique reappears
throughout the Canon, and it is worth naming now, because it is the single
most characteristic feature of how this book is set. A rubric interleaved
mid-sentence is a rubric that cannot be performed at the wrong moment.

\subsection*{Absolution, versicles, ascent}

\begin{appointed}{1018--1021 $\cdot$ \latincue{Misereatur}, \latincue{Indulgentiam}, versicles}
\textit{Focused source locator.} The appointed unit is identified by this incipit and marginal number; only phrases required by the analysis below are quoted.
\end{appointed}

Three things in this short stretch repay attention. The \latincue{Misereatur}
is said twice with changed pronouns: the ministers pray for the priest, and
the priest then prays for the ministers. Neither utterance is the sacramental
absolution of Penance, and the ministers do not absolve the priest. The
\latincue{Indulgentiam} likewise has the form of a petition:
\latincue{tribuat} asks that the almighty and merciful Lord grant indulgence,
absolution, and remission. Its first-person plural
\latincue{peccatorum nostrorum} includes priest and ministers in the same
prayer, while the priest's sign of the cross gives the petition the ritual
form of a blessing. These distinctions matter even though the general rubrics
name the whole block the confession ``with the absolution.'' And the versicles
are not a devotional filler: \latincue{Deus, tu conversus vivificabis
nos} is a request that God turn first, and \latincue{Ostende nobis, Domine,
misericordiam tuam. Et salutare tuum da nobis} asks for the mercy and the
\latincue{salutare} --- the saving thing, or the Saviour --- that the Mass is
about to give.

The priest then says \latincue{Oremus} aloud, goes up to the altar, and says
two prayers in a low voice.

\begin{appointed}{1022--1023 $\cdot$ \latincue{Aufer a nobis} and \latincue{Oramus te}}
\textit{Focused source locator.} The appointed unit is identified by this incipit and marginal number; only phrases required by the analysis below are quoted.
\end{appointed}

\begin{englishwitness}{xviii}
``Take away from us our iniquities, we beseech thee, O Lord, that we may be
worthy to enter with pure minds into the Holy of holies. Thro'. Amen.'' And:
``We beseech thee O Lord by the merits of thy saints, whose relics are here,
and of all the saints: that thou wouldst vouchsafe to forgive me all my sins.
Amen.''
\end{englishwitness}

\latincue{Sancta sanctorum} is a deliberate double meaning, and the 1861
translator caught it by capitalising: it is the Holy of Holies of the temple,
and it is also the Roman chapel of that name at the Lateran. Fortescue notes
that these two prayers were said in the medieval rite of the papal chapel, and
that the words \latincue{sancta sanctorum} and \latincue{quorum reliquiae hic
sunt} had a particular local aptness there, where the relics really were
underfoot.\footnote{Adrian Fortescue, \work{The Mass: A Study of the Roman
Liturgy} (London and New York: Longmans, Green and Co., 1922), p.~227 and its
first footnote, read at the page image.} That is a historical observation, not
an argument that the prayers mean less elsewhere: in the rite's received
setting, the second prayer itself identifies the saints' relics as present
there.

The kiss of the altar is printed inside the prayer, at the exact clause
\latincue{quorum reliquiae hic sunt}. The gesture and the relative clause point
at the same stone.

\begin{proposal}{Editorial reading: why the preparation is at the foot and not at the altar}
The 1962 book gives no explanation of why these prayers are said below the
step. The layout suggests one. Everything from \latincue{In nomine Patris} to
\latincue{Et cum spiritu tuo} is said at the foot; \latincue{Aufer a nobis} is
said while going up; \latincue{Oramus te} is said with the hands on the altar
and ends with a kiss. The three prayers are thus a graded approach, and the
psalm that opens them is a psalm about approach. Fortescue makes the same
inference from the same fact --- ``That all this is rather preparation than
part of the Mass itself is shown by its recital at the foot of the altar,
before the celebrant goes up to it'' --- and adds that a bishop does not put on
the maniple until after the confession. This is a source-grounded reading of
the ritual arrangement; it is not a claim about what any medieval compiler
intended.
\end{proposal}
